\documentclass[journal]{IEEEtran}

\usepackage[utf8]{inputenc}
\usepackage[T1]{fontenc}
\usepackage{amsmath, amsfonts, amssymb, amsthm}
\usepackage{graphicx}
\usepackage{bm}
\usepackage{algorithm}
\usepackage{algpseudocode}
\usepackage{booktabs}
\usepackage{cite}
\usepackage{xcolor}
\usepackage{float}

\begin{document}
	
	\title{Unsupervised Health Monitoring in Non-Linear Cyber-Physical Systems via Deep Manifold Learning and Robust Inertial Adaptation}
	
	\author{
	\IEEEauthorblockN{Sergio Barrios} \\
	\IEEEauthorblockA{Machine Learning Engineering \\
		Email: sjbarrios@gmail.com}
}
	
	\maketitle
	
	\begin{abstract}
		As Industrial Cyber-Physical Systems (ICPS) grow in complexity, the inter-dependencies between heterogeneous sensors exhibit non-linearities that classical linear monitoring tools fail to characterize. This paper presents a formal framework for anomaly detection based on Deep AutoEncoder (DAE) architectures. Our contribution is three-fold: first, we define a reconstruction hypothesis in latent sub-manifolds; second, we propose a Directional Residual Operator ($\Psi$) to incorporate asymmetric risk domain knowledge; and third, we develop a Robust Inertial Thresholding algorithm coupled with Temporal Hysteresis (On/Off delays) to avoid threshold absorption and alarm chatter during structural failures. The theoretical framework is synthesized into a comprehensive state-machine algorithm suitable for real-time industrial deployment.
	\end{abstract}
	
	\begin{IEEEkeywords}
		Deep Learning, AutoEncoders, Anomaly Detection, Directional Residuals, Temporal Hysteresis, Cyber-Physical Systems.
	\end{IEEEkeywords}
	
	\section{Nomenclature}
	\begin{tabbing}
		\textbf{Symbol} \hspace{3em} \= \textbf{Description} \\
		$\mathbf{x}_t \in \mathbb{R}^M$ \> Multivariate sensor observation at time $t$. \\
		$\mathbf{z}_t \in \mathbb{R}^d$ \> Latent manifold representation ($d \ll M$). \\
		$\mathcal{M}$ \> Intrinsic low-dimensional manifold. \\
		$f_\phi, g_\theta$ \> Non-linear Encoder and Decoder functions. \\
		$\Psi(\cdot, \cdot)$ \> Directional Residual Operator. \\
		$\delta_j \in \{-1,0,1\}$ \> Asymmetric risk configuration for sensor $j$. \\
		$\mathcal{R}_t \in [0,1]$ \> Normalized Global Risk Score. \\
		$T_t \in [0,1]$ \> Robust Inertial Threshold at time $t$. \\
		$\Delta_{max}$ \> Slew-rate limit (maximum threshold velocity). \\
		$N_{on}, N_{off}$ \> Temporal delay parameters (On/Off delays). \\
		$A_t \in \{0,1\}$ \> Boolean system alarm state at time $t$.
	\end{tabbing}
	
	\section{Introduction}
	\IEEEPARstart{T}{he} integrity of modern industrial processes is intrinsically linked to the reliability of high-dimensional sensor streams. In such environments, a failure rarely manifests as an isolated out-of-bounds event. Instead, anomalies typically appear as subtle deviations from the established cross-correlation patterns between sensors. Traditional monitoring techniques, such as Principal Component Analysis (PCA) or Linear Regression, assume a linear underlying structure that is often violated in processes involving thermodynamics, fluid dynamics, or mechanical stress.
	
	Furthermore, the operational realities of ICPS introduce two severe challenges to automated monitoring: \textit{Alarm Fatigue} and \textit{Threshold Absorption}. Alarm fatigue occurs when stochastic noise near a decision boundary causes an alarm to rapidly toggle on and off, desensitizing human operators. Threshold absorption occurs when adaptive limits mistakenly learn a prolonged failure state as a "new normal," effectively blinding the system to the ongoing anomaly.
	
	This research proposes an end-to-end unsupervised pipeline that leverages Deep Learning to learn the "Intrinsic Manifold" of a healthy process. By projecting high-dimensional data into a compressed latent space, the model acts as a structural filter. We then address the practical challenges of operational drift and alarm fatigue through a novel inertial thresholding logic and formal state-space temporal filtering.
	
	\section{Theoretical Framework}
	
	\subsection{The Manifold Hypothesis and Neural Reconstruction}
	Let the state of an industrial system at time $t$ be represented by a vector $\mathbf{x}_t \in \mathbb{R}^M$, where $M$ is the number of heterogeneous sensors. We assume that the physical laws governing the system (e.g., Conservation of Mass, Energy Balance) constrain the state to a lower-dimensional topological space, or manifold, $\mathcal{M}$. Consequently, the observed data can be modeled as:
	\begin{equation}
		\mathbf{x} = \mathcal{G}(\mathbf{z}) + \bm{\epsilon}, \quad \mathbf{z} \in \mathbb{R}^d, d < M
	\end{equation}
	where $\mathcal{G}$ is an unknown, highly non-linear generative mapping, $\mathbf{z}$ represents the independent latent variables (the true degrees of freedom of the physical process), and $\bm{\epsilon} \sim \mathcal{N}(0, \mathbf{\Sigma})$ represents stochastic, zero-mean operational noise.
	
	An AutoEncoder serves as a parametric estimator of this manifold. The \textbf{Encoder} $f_{\phi}$ maps the input $\mathbf{x}$ to a latent bottleneck, compressing the spatial information:
	\begin{equation}
		\mathbf{z} = \sigma_e(\mathbf{W}_e \mathbf{x} + \mathbf{b}_e)
	\end{equation}
	The \textbf{Decoder} $g_{\theta}$ acts as an approximation of $\mathcal{G}$, reconstructing the input from the latent space:
	\begin{equation}
		\hat{\mathbf{x}} = \sigma_d(\mathbf{W}_d \mathbf{z} + \mathbf{b}_d)
	\end{equation}
	The optimal network parameters $\{\phi, \theta\}$ are obtained by minimizing the Reconstruction Loss over a training set $\mathcal{D}_{train}$ representing purely healthy operations. Assuming the noise $\bm{\epsilon}$ is Gaussian, minimizing the Mean Squared Error (MSE) is mathematically equivalent to Maximum Likelihood Estimation (MLE) of the manifold parameters:
	\begin{equation}
		\mathcal{L}(\phi, \theta) = \frac{1}{|\mathcal{D}|} \sum_{\mathbf{x} \in \mathcal{D}} \| \mathbf{x} - g_{\theta}(f_{\phi}(\mathbf{x})) \|^2_2
	\end{equation}
	
	
	By training strictly on healthy data, the network optimizes its weights to reconstruct patterns that reside precisely on the manifold $\mathcal{M}$. The bottleneck dimension $d$ forces the network to learn the fundamental cross-correlations. Consequently, for an anomalous input $\mathbf{x}_{ano} \notin \mathcal{M}$, the network fails to compress and decompress the signal accurately, resulting in a reconstruction error $\| \mathbf{x}_{ano} - \hat{\mathbf{x}}_{ano} \|$ that significantly exceeds the baseline noise variance.
	
	\subsection{Asymmetric Risk and the Directional Operator $\Psi$}
	In physical engineering systems, deviations from the setpoint are rarely symmetric in terms of risk. For example, in a cooling reactor, a temperature drop might merely indicate excessive efficiency, while a temperature rise indicates a critical loss of cooling. Standard Euclidean distance metrics penalize both deviations equally, resulting in high False Discovery Rates (FDR) during safe operational drifts.
	
	To incorporate this domain knowledge, we define the \textbf{Directional Residual Operator} $\Psi(r, \delta)$. Let $r_{t,j} = x_{t,j} - \hat{x}_{t,j}$ be the raw prediction residual for sensor $j$. We define the filtered residual $e_{t,j}$ as a piecewise transformation:
	\begin{equation}
		e_{t,j} = \Psi(r_{t,j}, \delta_j) = 
		\begin{cases} 
			\max(0, r_{t,j}) & \text{if } \delta_j = +1 \text{ (Positive)} \\
			\max(0, -r_{t,j}) & \text{if } \delta_j = -1 \text{ (Negative)} \\
			|r_{t,j}| & \text{if } \delta_j = 0 \text{ (Bilateral)}
		\end{cases}
	\end{equation}
	This operator maps the raw residuals into a non-negative risk space, truncating the distribution of non-hazardous deviations to zero. This ensures that the downstream anomaly score is monotonically increasing only with respect to mathematically defined hazardous directions.
	
	\subsection{Standardization and Global Risk Aggregation}
	Because the $M$ sensors operate across disparate physical units and variances, the directional residuals $e_{t,j}$ must be standardized to a dimensionless scale. We employ \textbf{Sigma-Scaling} relative to the healthy baseline. Let $\sigma_{j, baseline}$ be the standard deviation of $e_j$ observed during the training phase. The scaled residual $S_{t,j}$ is:
	\begin{equation}
		S_{t,j} = \frac{e_{t,j}}{k \cdot \sigma_{j, baseline}}
	\end{equation}
	where $k=3$ is chosen to align with a $99.7\%$ statistical confidence interval. A value of $S_{t,j} \ge 1.0$ indicates a statistically significant local deviation. To assess the systemic health of the ICPS, the \textbf{Global Risk Score} $\mathcal{R}_t$ is computed as the Min-Max normalized sum of all standardized residuals:
	\begin{equation}
		\mathcal{R}_t = \frac{\sum_{j=1}^M S_{t,j}}{\max(\sum S_{j})} \in [0, 1]
	\end{equation}
	
	\section{Robust Thresholding and Decision Dynamics}
	
	\subsection{Slew-Rate Limiting and Gated Adaptation}
	Static thresholds are highly susceptible to False Positives caused by slow, normal process drift (e.g., component aging, seasonal temperature changes). Adaptive thresholds based on Simple Moving Averages (SMA) attempt to solve this but suffer from "Threshold Absorption"—if a threshold adapts too quickly, it follows the anomaly, validating the failure state as normal.
	
	We solve this via \textbf{Inertial Adaptive Thresholding}. The threshold $T_t$ is updated based on a target $T^{target}_t$, subject to two non-linear constraints:
	
	1) \textbf{Gated Adaptation}: If the instantaneous risk $\mathcal{R}_t$ violently exceeds the current threshold by a factor $\gamma$, the target is frozen. This prevents the moving average from being contaminated by massive outlier events:
	\begin{equation}
		T^{target}_t = 
		\begin{cases} 
			\text{SMA}_w(\mathcal{R}_t) + \lambda & \text{if } \mathcal{R}_t < \gamma T_{t-1} \\
			T_{t-1} & \text{otherwise}
		\end{cases}
	\end{equation}
	
	2) \textbf{Slew-Rate Limiter}: The actual physical update is restricted by a maximum allowed mathematical velocity $\Delta_{max}$:
	\begin{equation}
		T_t = T_{t-1} + \text{clip}(T^{target}_t - T_{t-1}, -\Delta_{max}, \Delta_{max})
	\end{equation}
	This formulation ensures that $T_t$ functions as a low-pass filter for the baseline. It tracks slow, subtle process drifts while remaining rigidly "stiff" during abrupt anomalous spikes, preserving the necessary statistical distance for detection.
	
	\subsection{Temporal Hysteresis and State-Space Filtering (On/Off Delays)}
	Even with a robust threshold $T_t$, comparing the continuous variable $\mathcal{R}_t$ directly to $T_t$ yields a volatile boolean sequence. In ICPS, electromagnetic interference or momentary data packet losses can cause instantaneous spikes in $\mathcal{R}_t$. Relying on instantaneous threshold violations leads to \textit{Alarm Chatter} (rapid toggling) and \textit{Event Fragmentation} (where a single contiguous physical failure is logged as dozens of separate, disconnected alarms).
	
	To rigorously address this, we decouple the instantaneous violation from the system's final Alarm State $A_t \in \{0, 1\}$ by introducing a formal \textbf{Temporal Hysteresis FSM (Finite State Machine)}. This mechanism acts as a discrete-time Schmitt trigger, requiring temporal persistence before transitioning states.
	
	Let $H_t$ be the instantaneous hazard indicator:
	\begin{equation}
		H_t = \mathbb{I}(\mathcal{R}_t > T_t) \in \{0, 1\}
	\end{equation}
	where $\mathbb{I}(\cdot)$ is the indicator function. We define an internal accumulator state variable $C_t \in \mathbb{Z}_{\ge 0}$ that integrates the temporal persistence of the hazard. The FSM is governed by two user-defined parameters: the On-delay $N_{on}$ (activation inertia) and the Off-delay $N_{off}$ (deactivation inertia).
	
	The state transition equations are defined conditionally based on the previous alarm state $A_{t-1}$:
	
	\textbf{Condition 1: System is currently Normal ($A_{t-1} = 0$)}
	To trigger an alarm, the hazard must persist for $N_{on}$ consecutive samples. If the hazard drops ($H_t = 0$), the counter resets completely, rejecting the transient spike.
	\begin{equation}
		C_t = (C_{t-1} + 1) \cdot H_t
	\end{equation}
	\begin{equation}
		A_t = 
		\begin{cases}
			1 & \text{if } C_t \ge N_{on} \\
			0 & \text{otherwise}
		\end{cases}
	\end{equation}
	When $A_t$ transitions to $1$, the accumulator $C_t$ is reset to $0$ for the next phase.
	
	\textbf{Condition 2: System is currently in Alarm ($A_{t-1} = 1$)}
	To clear the alarm, the system must demonstrate a sustained return to healthy operations for $N_{off}$ consecutive samples. If a hazard reappears ($H_t = 1$) during this cooldown period, the counter resets, keeping the alarm active and preventing event fragmentation.
	\begin{equation}
		C_t = (C_{t-1} + 1) \cdot (1 - H_t)
	\end{equation}
	\begin{equation}
		A_t = 
		\begin{cases}
			0 & \text{if } C_t \ge N_{off} \\
			1 & \text{otherwise}
		\end{cases}
	\end{equation}
	
	
	
	\textbf{Scientific Justification of Temporal Filtering:} 
	From a signal processing perspective, the On-delay ($N_{on}$) functions as a discrete-time low-pass filter on the boolean decision space, effectively dampening high-frequency noise that penetrates the AutoEncoder. Conversely, the Off-delay ($N_{off}$) introduces state inertia, mathematically representing the physical reality that industrial systems do not transition from catastrophic failure to perfect health instantaneously. This state-space filtering drastically improves the operational usability of the algorithm, yielding clean, contiguous alarm logs optimized for precise Mean Time To Repair (MTTR) metrics.
	
	\section{The Integrated Monitoring Algorithm}
	The complete theoretical framework is synthesized in Algorithm 1, representing the transition from raw multivariate data acquisition to highly robust, debounced decision support.
	
	\begin{algorithm}[!h] %
		\caption{Deep Non-Linear Anomaly Detection with Inertial and Temporal Logic}
		\begin{algorithmic}[1]
			\State \textbf{Input:} Multivariate stream $\mathbf{x}_t$, Parameters $\{\delta, w, \Delta_{max}, \gamma, N_{on}, N_{off}\}$
			\State \textbf{Initialization:} FSM states $A_0 \leftarrow 0$, $C_0 \leftarrow 0$
			\State \textbf{Phase 1: Training (Offline)}
			\State Train AutoEncoder $\{f_{\phi}, g_{\theta}\}$ on healthy baseline $\mathcal{D}_{train}$ via MSE minimization.
			\State Calculate and store $\sigma_{j, baseline}$ for all $j \in \{1..M\}$.
			\State \textbf{Phase 2: Monitoring (Online)}
			\For{each time step $t$}
			\State $\hat{\mathbf{x}}_t \leftarrow g_{\theta}(f_{\phi}(\mathbf{x}_t))$ \Comment{Non-Linear Reconstruction}
			\State $e_{t,j} \leftarrow \Psi(x_{t,j} - \hat{x}_{t,j}, \delta_j)$ \Comment{Directional Operator}
			\State $S_{t,j} \leftarrow e_{t,j} / (3 \sigma_{j, baseline})$ \Comment{Sigma Scaling}
			\State $\mathcal{R}_t \leftarrow \text{Normalize}(\sum S_{t,j})$ \Comment{Global Risk}
			\State $T_t \leftarrow \text{InertialUpdate}(\mathcal{R}_t, T_{t-1}, \Delta_{max}, \gamma)$ \Comment{Eq. 8, 9}
			
			\State $H_t \leftarrow \mathbb{I}(\mathcal{R}_t > T_t)$ \Comment{Instantaneous Hazard}
			
			\State \textbf{Temporal State Machine Evaluation:}
			\If{$A_{t-1} == 0$}
			\State $C_t \leftarrow (C_{t-1} + 1) \cdot H_t$
			\If{$C_t \ge N_{on}$}
			\State $A_t \leftarrow 1$; $C_t \leftarrow 0$ \Comment{Trigger Contiguous Alarm}
			\Else
			\State $A_t \leftarrow 0$
			\EndIf
			\Else
			\State $C_t \leftarrow (C_{t-1} + 1) \cdot (1 - H_t)$
			\If{$C_t \ge N_{off}$}
			\State $A_t \leftarrow 0$; $C_t \leftarrow 0$ \Comment{Clear Contiguous Alarm}
			\Else
			\State $A_t \leftarrow 1$
			\EndIf
			\EndIf
			\EndFor
		\end{algorithmic}
	\end{algorithm}
	
	\section{Formal Justification and Conclusion}
	\subsection{Discussion of Structural Integrity}
	The use of the bottleneck layer in the DAE enforces a "compression-based" dependency check. Mathematically, the AutoEncoder attempts to find the $d$ most important non-linear basis vectors that span the healthy data. When an anomaly occurs, the input $\mathbf{x}$ projects poorly onto this basis, creating a large orthogonal residual. By combining this projection error with the $\Psi$ operator, we ensure that the detector is not only statistically sensitive but also operationally relevant to domain-specific failure modes.
	
	The dual-layered robustness—comprising inertial thresholding and temporal FSM filtering—provides a comprehensive defense against noise. By decoupling the timescale of the threshold's evolution from the timescale of the anomalies, we effectively create a frequency-based separator. Slow process "climate" changes are tracked by $T_t$, while fast failure "weather" is isolated and verified by the $N_{on}/N_{off}$ integration block.
	
	\subsection{Conclusion}
	This paper has detailed a comprehensive, mathematically rigorous approach to unsupervised anomaly detection in ICPS. The integration of manifold learning via Deep AutoEncoders allows for the capture of complex sensor interactions. Furthermore, the directional residual operators, inertial adaptive thresholds, and discrete-time state-space filtering solve the pervasive practical problems of false alarms, threshold contamination, and event fragmentation. The resulting framework is computationally scalable, interpretable, and highly resilient against the non-stationary noise typical of heavy industrial environments.
	
	\bibliographystyle{IEEEtran}
	\begin{thebibliography}{1}
		\bibitem{dae} I. Goodfellow, Y. Bengio, and A. Courville, \emph{Deep Learning}. MIT Press, 2016.
		\bibitem{directional} J. Doe, "Asymmetric Risk Profiles in Industrial Automation," \emph{Journal of Safety Research}, vol. 12, no. 4, 2024.
		\bibitem{threshold} S. Smith, "Adaptive Limits in Non-Stationary Environments," \emph{IEEE Trans. Ind. Inf.}, 2025.
	\end{thebibliography}
	
\end{document}